\section{The readout hypothesis}\label{sec:readout}

Sections~\ref{sec:kernel}--\ref{sec:generation} use neutral names.  This section states the physical hypothesis that motivated them, what in
the structure supports it, and what it would need in order to be more than a reading.

\subsection{Statement}

\begin{quote}
\emph{Readout hypothesis.}  Matter, light and a tachyon sector are one object read out in three frames.  The matter and the tachyon reading
differ, ideally, only in the direction of time; light, as the frame-invariant reading, has no time orientation.
\end{quote}

On the kernel the three readings are the positive, the negative and the null sector.  The hypothesis therefore makes two claims that can
be checked against the structure: the sectors must look alike in every respect that does not involve the direction of time, and the map
between the positive and the negative reading must contain time reversal.

\subsection{What the structure shows}

The first claim holds in a strong sense: the two sectors form one orbit (Proposition~\ref{prop:duality}), and every statement of
Sections~\ref{sec:kernel} and~\ref{sec:frames} holds for both alike.  For the second claim there are three points at which time reversal appears.
\begin{enumerate}
\item \emph{The involution of the compensator contains time reversal.}  By \eqref{eq:JB}, \(J_B=R_{\pi/2}\circ T\), and \(J_B\) maps the positive to the
negative sector, \(s(J_Bz)=-\overline{s(z)}\).  The duality quarter turn alone also exchanges the sectors, but reverses the second invariant
\(\operatorname{Im}s=2E\cdot B\); \(J_B\), which contains \(T\), keeps it.  On the state disk \(J_B\) is the reflection of the axis between the sectors,
and the null centre is fixed.
\item \emph{On the generation, time reversal exchanges the two halves.}  \(T\) maps the image of \(C_{\rm ch}\) to its kernel and reverses all charges,
\(TQ_{\rm em}T^{-1}=-Q_{\rm em}\), and it exchanges the two anomaly-free projectors (Propositions~\ref{prop:generation} and~\ref{prop:anomaly}).  Read through
the hypothesis, the matter and the tachyon frame are, on one generation, particle and antiparticle content of the same algebra.
\item \emph{On the null set of the norm the negative sector needs an imaginary parameter.}  On \(N(q)=-t^2+s=0\) one has \(t^2=s\).  For the positive
sector \(t=\pm\sqrt s\) is real; for the negative sector \(t\) must be imaginary.  This is the formal signature of a tachyon, whose mass
parameter is imaginary.
\end{enumerate}

\subsection{What the hypothesis does not yet do}

That antiparticles can be read as particles propagating backwards in time is the Feynman--St\"uckelberg interpretation~\cite{stueckelberg,feynman}.
Points~1 and~2 are compatible with it and do not go beyond it.  A tachyon, in contrast, has a space-like momentum, and nothing in the kernel
produces one: the negative sector consists of ordinary magnetic-dominated fields, and the imaginary \(t\) of point~3 is a parameter of the
kernel, not a momentum.  The step from ``time-reversed reading'' to ``space-like propagation'' is therefore the open core of the hypothesis.
It would need a dynamics in which the parameter \(t\) becomes a frequency or a mass, so that its imaginary value has kinematic consequences.

The hypothesis could fail in two ways.  If a dynamics is found in which the positive and the negative sector differ in more than the
direction of time, the hypothesis is false in its strong form.  And if the map between the sectors that such a dynamics singles out does not
contain time reversal, the three readings are not related as claimed.  Neither test can be made without a dynamics
(Section~\ref{sec:discussion}).
